\itemhead{习题2.2 · 第4题（教材p.49）}{XJ-B1-C02-HW-P049-Y0793}
\label{XJ-B1-C02-HW-P049-Y0793}
{\footnotesize 湘教版教材p.49 / PDF第56页 / 习题2.2 / 第4题}
\field{原题}{如图所示为二次函数 $y=ax^2+bx$ 的图象，若一元二次方程 $ax^2+bx+m=0$ 有实数根，求 $m$ 的最大值。
\bankfig{C02-HW049}{0.36}}
\field{参考答案}{图示原函数最低点纵坐标为 $-3$，故方程有实根等价于 $m\le3$，最大值为3。}
{\small \field{关联知识}{Q01 二次函数、零点、判别式与根系关系（教材帮p.65 / PDF74）；F04 函数三种表示法、图象与读数（教材帮p.85、86、93、94 / PDF94、95、102、103）}
\field{方法与易错点}{按已知零点、顶点或系数选择交点式、顶点式或一般式；根平方和用根系关系消元。 易错：遗漏 $a
e0$；把二重根算成两个不同零点；有两正根只查积正。}
\field{答案依据}{答案册相关题组：PDF第13、14页（册内页12、13）；按当前题库原题逐项重算，不按旧版题号机械对套。}}
